\clearpage
\begingroup
\fancyhead[L]{\small BEMERKUNGEN}
\renewcommand{\headrulewidth}{0pt}
\par\noindent\hspace*{\parindent}Im Einzelnen ist folgendes zu bemerken:

\par\noindent\hspace*{\parindent}Zu [I.]: Die Briefstellen sind nach den Originalen abgedruckt,
mit Ausnahme des ersten Briefes an \textsc{Olbers} vom 3. Januar 1806,
dessen Original sich auf der Pulkowaer Sternwarte befindet und von dem das
Archiv durch Vermittlung von Herrn C. \textsc{Schilling}-Bremen eine Copie besitzt.
Zum Brief von \textsc{Hansen} an \textsc{Gauss} vom 7. Februar 1843 ist zu bemerken,
dass nicht \(\dfrac{p''}{p_1''\mu}\), sondern das Mittel
\(\frac12\left(\dfrac{p''}{p_1''\mu}+\dfrac{r''}{r_1''\mu}\right)\)
die letzterwähnte Zahl \(9{,}791\,4946_n\) ergibt, also nicht
\(\dfrac{p''}{p_1''\mu}\) mit \(\dfrac{p'''}{p_1'''\mu}\) übereinstimmt,
sondern dies Mittel; ich verdanke diese Mittheilung Herrn J. \textsc{Kramer},
der das Beispiel nachgerechnet hat.

\par\noindent\hspace*{\parindent}Zu [II.]: Von der »Exposition d'une nouvelle méthode« etc.
befinden sich im Nachlass zwei unvollendete Entwürfe, von denen nur der erste
das mitabgedruckte Motto trägt. Die ersten drei »Sections« unterscheiden sich
nicht wesentlich in beiden Entwürfen und sind (Seite 439--464) nach dem zweiten
Entwurf abgedruckt. Die »Quatrième Section« dagegen ist in beiden Entwürfen
gänzlich verschieden und darum nach beiden abgedruckt; im ersten Entwurf bricht
sie inmitten eines Satzes ab. Bei den Manuscripten fand sich auf einem Zettel
folgende »L[\textsc{isting}] 1855 Aug. 18« unterzeichnete Bemerkung:
»Diese offenbar ganz druckfertige Abhandlung, welche sich in dem \((3+4)\)ten Bande
von \textsc{Crelles} Journal vorfand, scheint in's Jahr 1806 (spätestens) gesetzt
werden [zu] müssen, insofern die in der »Quatrième Section« enthaltene
Entwickelungsmethode nach \textsc{Gauss}' mündlichen Äusserungen früher als die von
\textsc{Fourier} gegebene vom Dec. 1806 niedergeschrieben war«, mit dem späteren Zusatz:
»\textls[120]{Jedoch} stellt die Stelle Art. 20 »pendant les années 1803 \ldots{} 1816«
die Zeit der Abfassung nach dem letztgenannten Jahre. Noch wäre nicht ohne Interesse,
dem Ursprung der zu Ende Art. 3 angeführten Constanten nachzuspüren.«

\par\noindent\hspace*{\parindent}Aus dem Briefwechsel geht aber hervor, dass jedenfalls der erste
Entwurf etwa 1815 zu setzen ist, der zweite, der den von \textsc{Listing} erwähnten
Gegenstand behandelt, also noch später; \textsc{Gauss} mündliche Äusserungen können sich
nur darauf beziehen, dass er diese Untersuchungen überhaupt schon gemacht,
eventuell an anderer Stelle aufgezeichnet hatte. Von befreundeter Seite erhielt
ich einige interessante Mittheilungen über den \textsc{Gauss}schen Convergenzbeweis,
von denen ich hier nur erwähnen kann, dass \textsc{Poissons} analoger Beweis bereits
1815 im Bulletin der Société philomatique p. 90 im Auszug erschienen ist.
Es wäre nicht unmöglich, dass \textsc{Gauss} diesen Auszug bei der Abfassung des zweiten
Entwurfs gekannt hat und vielleicht gerade dadurch zum Niederschreiben des von ihm
früher gefundenen Beweises veranlasst wurde. \textsc{Gauss} scheint, wie schon erwähnt,
die Ausarbeitung der »Exposition« in der Absicht unternommen zu haben, sie für den
von der Pariser Akademie auf die Störungen der Pallas ausgesetzten Preis
einzusenden, dessen Vertheilung schliesslich bis zum Jahr 1816 hinausgeschoben
worden war (vgl. auch den Brief an \textsc{Olbers} vom 8. Januar 1816, S. 427).
\endgroup
